\begin{lemma}
Let $U$ be a convex open subset of $\mathbb{R}^n$ containing the origin, let
$M\ge0$ and $\alpha\in(0,1]$, and let $F\colon\mathbb{R}^n\to\mathbb{C}$ be a
function of class $C(U)\cap C^1(U-\{0\})$ such that
\begin{enumerate}
\item[(i)] $|F(x)|\le M|x|^\alpha$ for $x\in U$,
\item[(ii)] $\|DF(x)\|\le M|x|^{\alpha-1}$ for $x\in U-\{0\}$,
\end{enumerate}
where $DF(x)$ denotes the real Fr\'echet derivative of $F$ at $x$, a
real-linear map $\mathbb{R}^n\to\mathbb{C}$, and $\|DF(x)\|$ its operator
norm.  Then
\[
  |F(x)-F(y)|\le5M|x-y|^\alpha\qquad\text{for all }x,y\in U .
\]
\end{lemma}
\begin{proof}
Let $x,y\in U$.

\emph{The degenerate case $x=y$.}  Then $|F(x)-F(y)|=0$, and since
$\alpha>0$ we have $|x-y|^\alpha=0^\alpha=0$, so the right-hand side is
$0$ as well and the inequality holds.  Assume from now on that $x\neq y$, so
that $|x-y|>0$.

Exactly one of the following two cases occurs.

\emph{Case 1: $|y|\le2|x-y|$.}  By the triangle inequality in $\mathbb{C}$,
\[
  |F(x)-F(y)|\le|F(x)|+|F(y)| .
\]
Now apply the unweighted two-term power-mean inequality
$(z_1+z_2)^{p}\le2^{p-1}(z_1^{p}+z_2^{p})$, valid for $z_1,z_2\ge0$ and
$p\ge1$ (mathlib's \texttt{NNReal.rpow\_add\_le\_mul\_rpow\_add\_rpow}), at the
exponent $p=1/\alpha$, which is $\ge1$ because $0<\alpha\le1$, and with
$z_1=|F(x)|\ge0$, $z_2=|F(y)|\ge0$.  Raising the result to the power $\alpha$,
which is monotone since $\alpha>0$, and using
$\bigl(t^{1/\alpha}\bigr)^{\alpha}=t$ for $t\ge0$ on the left and
$\bigl(2^{1/\alpha-1}\bigr)^{\alpha}=2^{1-\alpha}$ on the right, gives
\[
  |F(x)|+|F(y)|\le2^{1-\alpha}\Bigl(|F(x)|^{1/\alpha}+|F(y)|^{1/\alpha}\Bigr)^{\alpha}.
\]
By hypothesis (i), $|F(x)|\le M|x|^\alpha$ and $|F(y)|\le M|y|^\alpha$, all
four quantities being nonnegative.  Since $1/\alpha>0$, the map
$t\mapsto t^{1/\alpha}$ is nondecreasing on $[0,\infty)$, so
\[
  |F(x)|^{1/\alpha}\le\bigl(M|x|^\alpha\bigr)^{1/\alpha}
  =M^{1/\alpha}\bigl(|x|^{\alpha}\bigr)^{1/\alpha}=M^{1/\alpha}|x| ,
\]
and likewise $|F(y)|^{1/\alpha}\le M^{1/\alpha}|y|$, using
$(t^\alpha)^{1/\alpha}=t$ for $t\ge0$ and the multiplicativity of real powers
on nonnegative bases.  Since $\alpha>0$, the map $t\mapsto t^{\alpha}$ is
nondecreasing on $[0,\infty)$, whence
\[
  2^{1-\alpha}\Bigl(|F(x)|^{1/\alpha}+|F(y)|^{1/\alpha}\Bigr)^{\alpha}
  \le2^{1-\alpha}\Bigl(M^{1/\alpha}\bigl(|x|+|y|\bigr)\Bigr)^{\alpha}
  =2^{1-\alpha}M\bigl(|x|+|y|\bigr)^{\alpha},
\]
where we used $\bigl(M^{1/\alpha}\bigr)^{\alpha}=M$ (valid since $M\ge0$).

We now bound $|x|+|y|$.  By the triangle inequality $|x|\le|x-y|+|y|$, so
\[
  |x|+|y|\le|x-y|+2|y|\le|x-y|+4|x-y|=5|x-y| ,
\]
the second inequality being the case hypothesis $|y|\le2|x-y|$.  Hence,
again by monotonicity of $t\mapsto t^\alpha$ on $[0,\infty)$ and $M\ge0$,
\[
  |F(x)-F(y)|\le2^{1-\alpha}M\bigl(5|x-y|\bigr)^{\alpha}
  =2^{1-\alpha}5^{\alpha}M|x-y|^{\alpha}.
\]
Finally
\[
  2^{1-\alpha}5^{\alpha}=2\cdot\Bigl(\frac52\Bigr)^{\alpha}\le2\cdot\frac52=5 ,
\]
because $\frac52>1$ and $0<\alpha\le1$, so that
$t\mapsto(\frac52)^{t}$ is nondecreasing.  Since $M|x-y|^\alpha\ge0$, we
conclude
\[
  |F(x)-F(y)|\le5M|x-y|^{\alpha}.
\]

\emph{Case 2: $|y|>2|x-y|$.}  For $t\in[0,1]$ put
\[
  z(t)=tx+(1-t)y .
\]
Since $U$ is convex and $x,y\in U$, we have $z(t)\in U$ for every
$t\in[0,1]$.  Moreover $z(t)=y+t(x-y)$, so for $t\in[0,1]$,
\[
  |z(t)|\ge|y|-t\,|x-y|\ge|y|-|x-y|>2|x-y|-|x-y|=|x-y|>0 ,
\]
using $t\le1$ and the case hypothesis.  In particular $z(t)\neq0$, so
\[
  z(t)\in U-\{0\}\qquad\text{for every }t\in[0,1].
\]

Since $U$ is open and $\{0\}$ is closed, $U-\{0\}$ is open, and $F$ is of
class $C^1$ on it.  The segment $\{z(t):t\in[0,1]\}$ is contained in
$U-\{0\}$, so the map $t\mapsto F(z(t))$ is of class $C^1$ on $[0,1]$, with
derivative
\[
  \frac{d}{dt}F\bigl(z(t)\bigr)=DF\bigl(z(t)\bigr)[x-y]
\]
by the chain rule, $t\mapsto z(t)$ being affine with derivative the constant
vector $x-y$.  By the fundamental theorem of calculus, applied to the real and
imaginary parts of this $C^1$ function on $[0,1]$,
\[
  F(x)-F(y)=F\bigl(z(1)\bigr)-F\bigl(z(0)\bigr)
  =\int_0^1 DF\bigl(z(t)\bigr)[x-y]\,dt .
\]

We estimate the integrand.  For each $t\in[0,1]$, by the definition of the
operator norm,
\[
  \bigl|DF\bigl(z(t)\bigr)[x-y]\bigr|\le\bigl\|DF\bigl(z(t)\bigr)\bigr\|\,|x-y| ,
\]
and by hypothesis (ii), applicable since $z(t)\in U-\{0\}$,
\[
  \bigl\|DF\bigl(z(t)\bigr)\bigr\|\le M\,|z(t)|^{\alpha-1}.
\]
Now $\alpha-1\le0$, so $r\mapsto r^{\alpha-1}$ is nonincreasing on
$(0,\infty)$; since $|z(t)|\ge|x-y|>0$, this gives
\[
  |z(t)|^{\alpha-1}\le|x-y|^{\alpha-1}.
\]
As $M\ge0$, combining the last three displays yields, for every $t\in[0,1]$,
\[
  \bigl|DF\bigl(z(t)\bigr)[x-y]\bigr|\le M\,|x-y|^{\alpha-1}\,|x-y|
  =M\,|x-y|^{\alpha},
\]
the last step using $|x-y|>0$ and
$|x-y|^{\alpha-1}|x-y|=|x-y|^{(\alpha-1)+1}=|x-y|^{\alpha}$.

Therefore
\[
  |F(x)-F(y)|=\Bigl|\int_0^1 DF\bigl(z(t)\bigr)[x-y]\,dt\Bigr|
  \le\int_0^1\bigl|DF\bigl(z(t)\bigr)[x-y]\bigr|\,dt
  \le\int_0^1 M|x-y|^{\alpha}\,dt=M|x-y|^{\alpha}.
\]
Since $M|x-y|^{\alpha}\ge0$, we have $M|x-y|^\alpha\le5M|x-y|^\alpha$, so
\[
  |F(x)-F(y)|\le5M|x-y|^{\alpha}
\]
in Case 2 as well.

As Cases 1 and 2 are exhaustive, the estimate (53) holds for all $x,y\in U$.
\end{proof}
